% Work-MAT | 2. rocnik | pripravovana hodina 23
% Návaznost: uživatelské připravy.zip (hodiny 01–20), autorská učebnice na GitHubu.
% Realisticky.cz: didaktická inspirace v posloupnosti a typech úloh; zadání jsou vlastní.
\documentclass[11pt,a4paper]{article}
\input{../../sablony/pripravy/priprava.sty}
\nastavPripravu{2. ročník}{Exponenciální rovnice}{23}{Substituce se zápornými exponenty}
\begin{document}
\blok{Motivace a návaznost}
Připomeň pravidlo $a^{-x}=1/a^x$. Pokud se ve výrazu objeví $a^x$ i $a^{-x}$, jediná substituce stále stačí.
\poznamka{Orientační nabídka pro 45--55 minut: krátké zopakování a motivace, nový výklad, tři společné vzorové úlohy, samostatná práce. Vyber podle tempa třídy; zbývající zadání slouží jako rezerva.}
\blok{Výklad}
Při $t=a^x>0$ platí $a^{-x}=\frac{1}{t}$. Násobení rovnice výrazem $t$ je přípustné, protože $t>0$.
\blok{Společné příklady}
\begin{prikladReseniBox}[title={Příklad 1}]
Vyřeš, resp. urči: $2^x+2^{3-x}=6$.
\tcblower
\[\begin{aligned}t=2^x>0,\quad t+\frac{8}{t}=6 \\ t^2-6t+8=0\;\Longrightarrow\;(t-2)(t-4)=0 \\ K=\{1,2\}\end{aligned}\]
\end{prikladReseniBox}
\begin{prikladReseniBox}[title={Příklad 2}]
Vyřeš, resp. urči: $3^x+3^{2-x}=6$.
\tcblower
\[\begin{aligned}t=3^x>0,\quad t+\frac{9}{t}=6 \\ t^2-6t+9=0\;\Longrightarrow\;t=3 \\ x=1\end{aligned}\]
\end{prikladReseniBox}
\begin{prikladReseniBox}[title={Příklad 3}]
Vyřeš, resp. urči: $5^x+5^{3-x}=30$.
\tcblower
\[\begin{aligned}t=5^x>0,\quad t+\frac{125}{t}=30 \\ (t-5)(t-25)=0\;\Longrightarrow\;K=\{1,2\}\end{aligned}\]
\end{prikladReseniBox}
\poznamka{V každém kroku napiš přeměnu $a^{k-x}=a^k/t$. Upozorni, že logaritmy zatím nepotřebujeme.}
\blok{Samostatné procvičování}
\begin{zadaniReseniBox}[title={Samostatná práce 1}]
\begin{enumerate}[label=\alph*)]
\item $2^x+2^{3-x}=6$.
\item $3^x+3^{2-x}=6$.
\item $5^x+5^{3-x}=30$.
\item $2^x+2^{4-x}=8$.
\item $3^x+3^{4-x}=18$.
\item Vysvětli podmínku při násobení výrazem $t$..
\end{enumerate}
\tcblower
\begin{enumerate}[label=\alph*)]
\item $K=\{1,2\}$.
\item $x=1$.
\item $K=\{1,2\}$.
\item $t=4, x=2$.
\item $x=2$.
\item $t=a^x>0$, takže nenásobíme nulou..
\end{enumerate}
\end{zadaniReseniBox}
\blok{Dobrovolné rozšíření}
\blok{Úkol na doma}
Při početních úlohách napiš postup, zvolenou metodu a zkoušku v původním zadání. U soustav kontroluj obě rovnice; u nerovnic výsledek zakresli na osu.
\begin{zadaniReseniBox}[title={Domácí úkol}]
\begin{enumerate}[label=\alph*)]
\item $2^x+2^{3-x}=6$.
\item $3^x+3^{2-x}=6$.
\item $5^x+5^{3-x}=30$.
\item $2^x+2^{4-x}=8$.
\item $3^x+3^{4-x}=18$.
\item $2^x+2^{2-x}=5$.
\item Jak se vyjádří $4^{-x}$ pomocí $t=2^x$?.
\end{enumerate}
\tcblower
\begin{enumerate}[label=\alph*)]
\item $K=\{1,2\}$.
\item $x=1$.
\item $K=\{1,2\}$.
\item $x=2$.
\item $x=2$.
\item $t=1 nebo4, K=\{0,2\}$.
\item $t^{-2}$..
\end{enumerate}
\end{zadaniReseniBox}
\poznamka{Domácí úkol obsahuje 7 položek; předpokládaná domácí práce přibližně 15--30 minut. Po větších blocích jsou vědomě promíchané dříve probrané typy.}
\end{document}
